% Generated draft: requires mathematical and source review.
\documentclass[11pt]{article}
\usepackage[margin=25mm]{geometry}
\usepackage{fontspec}
\setmainfont{Latin Modern Roman}
\usepackage{amsmath,amssymb,mathtools,unicode-math}
\setmathfont{Latin Modern Math}
\usepackage{enumitem,needspace,xurl,hyperref}
\setlength{\parindent}{0pt}
\setlength{\parskip}{5pt}
\setlength{\emergencystretch}{3em}
\newcommand{\N}{\mathbb{N}}
\newcommand{\Z}{\mathbb{Z}}
\newcommand{\Q}{\mathbb{Q}}
\newcommand{\R}{\mathbb{R}}
\newcommand{\C}{\mathbb{C}}
\newcommand{\F}{\mathbb{F}}
\newcommand{\A}{\mathbb{A}}
\newcommand{\PP}{\mathbb{P}}
\newcommand{\E}{\mathbb{E}}
\newcommand{\eps}{\varepsilon}
\newcommand{\dd}{\,\mathrm{d}}
\begin{document}
\section*{UnsolvedMath ID: 800096}
\textbf{Source problem code:} catalogue entry
\subsection{Definitions and mathematical statement}
Let \(M\) be a smooth \(3\)-manifold and let \(r\) denote the regularity parameter intended by the source. Write \(\mathscr F^r(M)\) for the space of codimension-one \(C^r\) foliations on \(M\) that are oriented and co-oriented: the tangent planes to the leaves have a continuous orientation, and the normal line bundle has a continuous orientation. The admissible values of \(r\), the hypotheses on \(M\), boundary conventions, and the precise \(C^r\) topology on this space are not specified in the supplied record.

A leaf \(L\) is a connected maximal integral surface of a foliation. For an orientable surface \(L\), interpret its genus as
\[
 g(L)=\sup\{g(K):K\subseteq L\text{ is a compact connected subsurface}\}\in\{0,1,2,\ldots\}\cup\{\infty\}.
\]
This convention includes noncompact leaves; its intended use requires confirmation from the source.

Under the usual interpretation of genericity, a subset is residual if its complement is a countable union of nowhere dense sets. With the scope and topology above still to be specified, the requested assertion is that there exists a residual subset \(\mathcal G\subseteq\mathscr F^r(M)\) such that
\[
 \forall\mathcal F\in\mathcal G\quad
 \forall L\text{ a leaf of }\mathcal F,\qquad g(L)<2.
\]
Equivalently, every leaf of every foliation in this residual subset has genus \(0\) or \(1\). The residual subset concerns entire foliations, and the conclusion concerns every leaf of each such foliation. This is a formulation scheme: the supplied record does not determine the range of manifolds or regularities over which it is asserted.
\subsection{Short English statement}
Question 37 asks whether a \(C^r\)-generic oriented and co-oriented codimension-one foliation of a \(3\)-manifold has only leaves of genus \(0\) or \(1\). The regularity range, manifold hypotheses, topology for genericity, and genus convention for noncompact leaves require clarification.
\subsection{Sources}
\begin{itemize}
\item Local UnsolvedMath registry, permanent ID \(800096\), exact source code \texttt{catalogue entry}, title ``Foliations and Diffeomorphism Groups --- Question 37.'' The supplied record attributes the problem to Nariman and Yazdi (2025) and labels its status as open; this status is supplied curation, not an independently verified mathematical assessment.
\item Nariman and Yazdi, \emph{Problem List on Foliations and Diffeomorphism Groups} (2025), Question 37. Supplied source URL: \url{https://arxiv.org/abs/2503.18274}. The record reports access on 2026-10-16 and extraction by \texttt{arxiv\_latex}; these are supplied provenance details, not a claim of independent source access or verification.
\end{itemize}
\end{document}
